% Pregunta 305 (2019-febrero-OEP2018, n.º 16): modelo WS-PS; más competencia baja el margen (mu' < mu), la PS sube y el equilibrio pasa de E0 a E1 sobre la WS: sube el salario real y baja la tasa natural de paro. Autor: Víctor Gutiérrez Marcos
\documentclass[tikz,border=6pt]{standalone}
\input{tikz-tcee}
\begin{document}
\begin{tikzpicture}[tcee]
  \ejes{$u$}{$W/P$}{6}{5}
  % WS decreciente: W/P = 0,6 + 7/(u+1)
  \draw[curva1,domain=0.75:5.3,samples=60] plot (\x,{0.6+7/(\x+1)});
  \node[etiqueta,tceemorado,anchor=west] at (5.3,1.72) {$WS$};
  % PS inicial: W/P = A/(1+mu) = 2,4
  \draw[curva2] (0,2.4) -- (5.3,2.4) node[etiqueta,right] {$PS$};
  % PS tras la caída del margen: W/P = A/(1+mu') = 3,4
  \draw[curvanueva,desplazada] (0,3.4) -- (5.3,3.4) node[etiqueta,right,text=tceerojo] {$PS'$};
  \draw[flecha] (4.5,2.5) -- (4.5,3.3);
  \node[etiqueta,tceerojo,anchor=west,font=\footnotesize] at (4.55,2.9) {$\mu'<\mu$};
  % Equilibrios
  \proyeccion{(2.889,2.4)}{$u_n$}{$\frac{A}{1+\mu}$}
  \proyeccion{(1.5,3.4)}{$u_n'$}{$\frac{A}{1+\mu'}$}
  \node[punto,label={[etiqueta]above right:$E_0$}] at (2.889,2.4) {};
  \node[punto,fill=tceerojo,label={[etiqueta]above right:$E_1$}] at (1.5,3.4) {};
  \draw[flecha] (2.75,0.25) -- (1.65,0.25);
\end{tikzpicture}
\end{document}
